\documentclass[11pt]{article}

\usepackage[T1]{fontenc}
\usepackage{lmodern}
\usepackage[margin=1.05in]{geometry}
\usepackage{microtype}
\usepackage{amsmath,amssymb,amsthm,mathtools}
\usepackage{enumitem}
\usepackage{booktabs,tabularx,array,longtable}
\usepackage{listings}
\usepackage{xcolor}
\usepackage{fancyhdr}
\usepackage[colorlinks=true,linkcolor=blue!45!black,
            citecolor=blue!45!black,urlcolor=blue!55!black]{hyperref}
\hypersetup{
  pdftitle={Nonproperness Components in a Degree-Six Keller Family - Final Lean-Synchronized Edition AF-5},
  pdfsubject={Human proof and complete Lean 4 formalization of Theorems A and B},
  pdfkeywords={Keller map, Jacobian conjecture, nonproperness set, Lean 4}
}

\allowdisplaybreaks[2]
\setlength{\parindent}{0pt}
\setlength{\parskip}{0.55em}
\setlist[enumerate]{leftmargin=2.1em,itemsep=0.2em,topsep=0.3em}
\renewcommand{\thesection}{\Alph{section}}

\definecolor{codebg}{RGB}{246,248,250}
\definecolor{codeframe}{RGB}{210,215,220}
\lstset{
  basicstyle=\ttfamily\small,
  backgroundcolor=\color{codebg},
  frame=single,
  rulecolor=\color{codeframe},
  breaklines=true,
  columns=fullflexible,
  keepspaces=true,
  showstringspaces=false,
  aboveskip=0.7em,
  belowskip=0.7em
}

\newtheorem{theorem}{Theorem}[section]
\newtheorem{lemma}[theorem]{Lemma}
\newtheorem{corollary}[theorem]{Corollary}
\newtheorem*{theoremA}{Theorem A}
\newtheorem*{theoremB}{Theorem B}
\newtheorem*{formalresult}{Kernel-checked main result}
\newtheorem*{componentformula}{General component formula}

\newcommand{\C}{\mathbb{C}}
\newcommand{\A}{\mathbb{A}}
\newcommand{\Gm}{\mathbb{G}_{m}}
\newcommand{\Irr}{\operatorname{Irr}}
\newcommand{\Aut}{\operatorname{Aut}}

\pagestyle{fancy}
\fancyhf{}
\fancyhead[L]{\small Degree-six Keller inequivalence - AF-5 final}
\fancyhead[R]{\small Human proof / Theorems A and B checked in Lean}
\fancyfoot[C]{\thepage}
\renewcommand{\headrulewidth}{0.35pt}

\title{\textbf{Nonproperness Components in a Degree-Six Keller Family}\\[0.35em]
\large A human proof synchronized with the complete Lean 4 formalization\\[0.2em]
\normalsize Including the universal family and infinite inequivalent subfamily}
\author{}
\date{1 August 2026\\[0.25em]\small Protocol A--F edition, final revision AF-5}

\begin{document}

\maketitle

\begin{abstract}
For the degree-six subfamily of the polynomial Keller maps in Theorem 5.2 of
\emph{A Counterexample to the Jacobian Conjecture}, we determine the reduced
nonproperness set.  If the only nonzero deformation parameter is
$\eta_4=h\in\C[T]\setminus\{0\}$, then
\[
 S_{F_h}=D_h\cup V\bigl(p^6h(p)\bigr),
\]
where $D_h$ is one irreducible nonvertical hypersurface and contains no
vertical hyperplane.  Consequently
\[
 \#\Irr(S_{F_h})
 =2+\#\{\alpha\in\C^\times:h(\alpha)=0\},
\]
with roots counted without multiplicity.  This proves that CEX-004 and
CEX-006 are not left-right polynomially equivalent and yields infinitely many
pairwise inequivalent Keller counterexamples of generic degree six.

The accompanying Lean 4 project is pinned to Lean~4.32.2 and to the exact
mathlib revision recorded in Section~D.  It now proves the family-wide
Jacobian identity, the genuine function-field generic degree, the degree-one
obstruction for polynomial automorphisms, the reduced nonproperness equality
for every nonzero $h$, the universal component formula, and the explicit
infinite pairwise inequivalent family.  The parameter-free declarations
\texttt{theoremA} and \texttt{theoremB} assemble these facts exactly.

The final pinned replay completed successfully with 3507 build jobs.  Its
enforced audit parsed 92 unique \texttt{\#print axioms} reports and accepted
only \texttt{propext}, \texttt{Classical.choice}, and \texttt{Quot.sound}; the
project source scan found no project-local \texttt{axiom}, \texttt{sorry}, or
\texttt{admit}.  The release manifest and archive are verified separately.
All nonproperness and zero-locus equalities are equalities of underlying
reduced sets; no scheme-theoretic multiplicity statement is asserted.
\end{abstract}

\tableofcontents

\section{Theorem intake}

\subsection{Exact objects, assumptions, and equality notion}

For a polynomial map $G:\C^n\to\C^n$, write
\[
 S_G=
 \left\{b\in\C^n:
 \begin{array}{c}
 \text{there exist }u_k\in\C^n\text{ with }\lVert u_k\rVert\to\infty,\\
 \text{and }G(u_k)\to b
 \end{array}
 \right\}
\]
for its nonproperness set.  Throughout, $S_G$ denotes its underlying reduced
algebraic set.  Thus every equality of nonproperness loci below is equality of
sets, and every irreducible-component count ignores scheme multiplicities.

The base field is $\C$.  In the family $F_h$ we assume
$h\in\C[T]\setminus\{0\}$.  Left-right polynomial equivalence means
\[
 G=L\circ F\circ R,
 \qquad L,R\in\Aut_{\mathrm{poly}}(\C^3).
\]
The proof uses classical polynomial algebra, Euclidean topology on complex
affine space, and standard algebraic geometry for constructible images,
dimension, Zariski closure, and irreducibility.

The intermediate formal relation used by Lean is slightly more explicit.  It
requires Euclidean homeomorphisms $L$ and $R$, preservation of escaping
sequences by $R$, a target affine-Zariski homeomorphism $L_z$ compatible with
$L$ on points, and the identity $G=L\circ F\circ R$.  This relation is named
\texttt{AlgebraicEscapeLeftRightEquivalent}.  The completed
\texttt{PolynomialAutomorphism} module represents polynomial forward and
inverse coordinate triples with exact two-sided inverse laws, derives their
Euclidean and affine-Zariski homeomorphisms, proves escape preservation, and
shows that ordinary polynomial left-right equivalence supplies all data in the
intermediate interface.

\subsection{Claims and proof status}

\begin{center}
\begin{tabularx}{\textwidth}{>{\raggedright\arraybackslash\bfseries}p{0.20\textwidth}X >{\raggedright\arraybackslash}p{0.18\textwidth}}
\toprule
Item & Exact status & Level \\
\midrule
Human mathematics & All displayed mathematical theorems in this document are
proved from the stated hypotheses, with the standard geometric results
identified where used. & H1 \\
Universal Lean family theorem & For every nonzero $h$, Lean proves constant
Jacobian $-2$, function-field generic degree six, nonautomorphism, the exact
reduced nonproperness equality, hypersurface and dimension-two finite geometry,
and the actual component formula. & Checked, unconditional \\
Formal Theorem A & The parameter-free declaration \texttt{theoremA} proves the
complete degree-six Keller-counterexample certificates, counts $3$ and $2$,
and polynomial left-right inequivalence. & Checked, unconditional \\
Formal Theorem B & Lean proves the product family has count $m+2$, an infinite
range, and pairwise polynomial left-right inequivalence on that same range. &
Checked, unconditional \\
Polynomial left-right bridge & Polynomial coordinate automorphisms with
polynomial inverses induce the required Euclidean and affine-Zariski
homeomorphisms and preserve component counts. & Checked, unconditional \\
Equality notion & Reduced algebraic sets only; no scheme-theoretic
multiplicity is claimed. & Explicit \\
Trust boundary & No project-local \texttt{axiom}, \texttt{sorry}, or
\texttt{admit}; audited declarations use only \texttt{propext},
\texttt{Classical.choice}, and \texttt{Quot.sound}. & Audited \\
\bottomrule
\end{tabularx}
\end{center}

\begin{theoremA}
The degree-six Keller counterexamples CEX-004 and CEX-006 are not left-right
polynomially equivalent.  Their reduced nonproperness sets have respectively
three and two irreducible components:
\[
 \#\Irr(S_{\mathrm{CEX\text{-}004}})=3,
 \qquad
 \#\Irr(S_{\mathrm{CEX\text{-}006}})=2.
\]
Here left-right polynomial equivalence means the notion fixed above.
\end{theoremA}

\begin{theoremB}
There are infinitely many pairwise left-right polynomially inequivalent Keller
counterexamples $\C^3\to\C^3$, all having generic degree six.
\end{theoremB}

\begin{formalresult}
In the pinned Lean environment, for every nonzero $h\in\C[T]$,
\[
 \operatorname{NonpropernessSet}(F_h)
 =D_h\cup V\bigl(p^6h(p)\bigr),
 \qquad
 \#\Irr(S_{F_h})
 =2+\#\{\alpha\in\C^\times:h(\alpha)=0\}.
\]
The polynomial model has Jacobian determinant $-2$, genuine function-field
generic degree six, and no polynomial inverse.  The parameter-free Lean
\texttt{theoremA} gives the two explicit counts and their polynomial
left-right inequivalence; \texttt{theoremB} gives the infinite, pairwise
inequivalent product family with count $m+2$.
\end{formalresult}

The universal component formula is
\[
 \#\Irr(S_{F_h})
 =2+\#\{\alpha\in\C^\times:h(\alpha)=0\},
\]
where distinct roots are counted without multiplicity.  The cases of constant
$h$, a root at zero, repeated roots, and the empty product in the infinite
family are included in the theorem as stated.

\section{Proof route}

\subsection{Shape and mathematical material}

The central claim is an equality of reduced algebraic sets.  The human proof
uses double inclusion, followed by an irreducible-component count.  It combines
exact polynomial identities, a rational parametrization, elementary properness
and covering arguments, and irreducibility of an image closure.  The final
polynomial inequivalence statement is obtained from invariance of the
component count.

The completed Lean proof follows the universal route described by the human
argument, with every topological transition made explicit.  It reifies $F_h$
as a polynomial coordinate map and computes the formal Jacobian determinant.
It defines generic degree as an actual intermediate-field finrank and obtains
degree six through the inverse generator and a rational-function field model.
For nonproperness, it combines Keller local homeomorphisms, local properness,
covering-space monodromy, dense good targets with exactly six simple inverse
roots, and an \texttt{encard} bound of at most three on the leading locus.
Exact eliminants and mathlib's genuine \texttt{irreducibleComponents} then
close the reduced-set equality and component formula.  No numerical or
external CAS result is imported as a hypothesis.

\subsection{Proof obligations and their closure}

\begin{enumerate}
 \item reify the point map as \texttt{FhPolynomial} and prove its exact
       evaluation bridge;
 \item compute the family-wide formal Jacobian determinant as the constant
       polynomial $-2$;
 \item define function-field generic degree as a finrank, prove it is six for
       nonzero $h$, and prove polynomial automorphisms have degree one;
 \item derive the inverse-root equation and reconstruction identities;
 \item prove the finite critical image and its closure are irreducible,
       hypersurface, dimension two, dominant, and nonvertical;
 \item prove local fiber constancy away from nonproperness using proper
       restricted Keller coverings and monodromy;
 \item meet each such neighborhood with the dense good-target set and transport
       the exact six-point fiber count;
 \item prove every leading-locus fiber is finite with \texttt{encard} at most
       three, forcing fiber loss and nonproperness;
 \item combine this reverse inclusion with universal asymptotic exclusion and
       exact eliminants to obtain the reduced-set equality;
 \item identify and count the genuine irreducible components, with repeated
       roots discarded and the zero root handled once; and
 \item specialize to the explicit product family and transport component count
       through polynomial left-right equivalence.
\end{enumerate}

Items 1--11 are kernel-checked for every nonzero deformation polynomial.  The
two parameter-free paper theorems are direct assemblies of these endpoints.

\subsection{Representation and library audit}

Euclidean complex affine space is represented by
\texttt{C3 := Fin 3 -> Complex}.  Its affine-Zariski topology is constructed
on the separate type synonym \texttt{ZariskiC3}, preventing accidental mixing
of analytic and algebraic topologies.  The invariant
\texttt{algebraicComponentCount} is defined as the cardinality of mathlib's
actual \texttt{irreducibleComponents} of the corresponding Zariski subspace.

The finite component is the Zariski closure of a structured critical image.
The source of the critical map is the principal open $s\ne0$ in affine
$2$-space.  A reusable principal-open regular-map layer proves the critical
map Zariski-continuous and its image irreducible.  The universal eliminant
identifies its closure with one nonzero-polynomial zero locus, proves
nonverticality and dominance, and supports the exact topological Krull
dimension computation.  The analytic modules connect the same locus to the
sequence-defined Euclidean nonproperness set.  Section~D records the exact
module boundaries and endpoint declarations.

\section{Human proof}

\subsection{The degree-six subfamily}

Set
\[
 A=1+xy,
 \qquad
 B=A^2z+y^2(4+3xy),
 \qquad
 P=AB,
\]
and
\[
 Q=y+3xB,
 \qquad
 R=2x-3x^2y-x^3z.
\]
Fix $h\in\C[T]\setminus\{0\}$.  In Theorem 5.2 of \cite{source}, take
$\eta_4=h$ and $\eta_j=0$ for $j\ne4$.  The resulting map is
\[
 F_h=(P,Q_h,R_h),
\]
where
\[
 Q_h=Q+A^2x^4B^6h(P),
 \qquad
 R_h=R-\frac23x^6B^6h(P).
\]

On the dense chart $A\ne0$, put $s=x/A$.  If
$(p,q,r)=F_h(x,y,z)$, define
\[
 \phi_h(p,s)=3ps+p^6h(p)s^4,
 \qquad
 \theta_h(p,s)=ps^3+\frac23p^6h(p)s^6.
\]
Then
\[
 q=y+\phi_h(p,s),
 \qquad
 r=2s-ys^2-\theta_h(p,s).
\]
Eliminating $y=q-\phi_h(p,s)$ gives
\begin{equation}
 \boxed{
 \Omega_{p,q,r}(s)
 =a_h(p)s^6+2ps^3-qs^2+2s-r=0,
 \qquad
 a_h(p)=\frac13p^6h(p).}
 \label{eq:omega}
\end{equation}
Moreover,
\begin{equation}
 \frac{\partial\Omega}{\partial s}
 =2\bigl(1-s(q-\phi_h(p,s))\bigr).
 \label{eq:derivative}
\end{equation}
Writing
\[
 D=1-s(q-\phi_h(p,s))=1-sy,
\]
a finite root is simple exactly when $D\ne0$.  For such a root, the
corresponding source point is
\begin{equation}
 y=q-\phi_h(p,s),
 \qquad
 x=\frac{s}{D},
 \qquad
 z=pD^3-y^2(4-sy)D.
 \label{eq:reconstruction}
\end{equation}
Indeed, these expressions give $A=D^{-1}$ and $B=pD$, and hence
$F_h(x,y,z)=(p,q,r)$.

\begin{lemma}[Keller property and generic degree]
\label{lem:keller-degree}
For every nonzero $h$, the map $F_h$ satisfies
\[
 \det JF_h=-2,
 \qquad
 \mu(F_h)=6.
\]
\end{lemma}

\begin{proof}
First observe that
\[
 \partial_s\theta_h=s^2\partial_s\phi_h.
\]
In the coordinates $(P,y,s)$, this gives
\[
 \det\frac{\partial(P,Q_h,R_h)}{\partial(P,y,s)}
 =2-2ys-\partial_s\theta_h+s^2\partial_s\phi_h
 =2(1-ys)=\frac{2}{A}.
\]
On the other hand,
\[
 \det\frac{\partial(P,y,s)}{\partial(x,y,z)}=-A.
\]
For completeness, $P_z=A^3$ and $s_x=A^{-2}$ with $y$ fixed, so expansion
along the $y$-row gives the displayed determinant.  Thus $\det JF_h=-2$ on
$A\ne0$.  Both sides are polynomial functions, so the identity holds on all
of $\C^3$.

Now let $K=\C(p,q)$.  Equation \eqref{eq:omega} has the form
\[
 f_{p,q}(s)-r=0,
 \qquad
 f_{p,q}(S)=a_h(p)S^6+2pS^3-qS^2+2S.
\]
Because $h\ne0$, the element $a_h(p)$ is nonzero in $K$, so $f_{p,q}$ has
degree six.  The polynomial $T-f_{p,q}(S)$ generates a prime ideal of
$K[S,T]$, since the quotient by that ideal is $K[S]$.  It is therefore
irreducible in $K[S,T]$, and Gauss's lemma makes $f_{p,q}(S)-r$ irreducible
in $K(r)[S]$.  Finally, \eqref{eq:reconstruction} gives
\[
 \C(x,y,z)=\C(p,q,r)(s).
\]
Consequently the induced function-field extension has degree six, so
$\mu(F_h)=6$.
\end{proof}

\subsection{Elementary facts about nonproperness}

\begin{lemma}[Closedness]
\label{lem:closed}
For every polynomial map $G:\C^n\to\C^n$, the set $S_G$ is closed in the
Euclidean topology.
\end{lemma}

\begin{proof}
Suppose $b\notin S_G$.  Then there are $R>0$ and $\varepsilon>0$ such that
\[
 \lVert u\rVert>R
 \quad\Longrightarrow\quad
 \lVert G(u)-b\rVert\ge\varepsilon.
\]
Otherwise one could choose $u_k$ with $\lVert u_k\rVert>k$ and
$\lVert G(u_k)-b\rVert<1/k$, contradicting $b\notin S_G$.  Every point
within $\varepsilon/2$ of $b$ obeys the same conclusion with
$\varepsilon/2$.  Thus the complement of $S_G$ is open.
\end{proof}

\begin{lemma}[Fiber loss forces nonproperness]
\label{lem:fiber-loss}
Let $G:\C^n\to\C^n$ be a Keller map of generic degree $d$.  If
$b\notin S_G$, then
\[
 \#G^{-1}(b)=d.
\]
Equivalently, any target whose fiber contains fewer than $d$ points lies in
$S_G$.
\end{lemma}

\begin{proof}
If $b\notin S_G$, the argument in Lemma \ref{lem:closed} gives a connected
neighborhood $U$ of $b$ over which $G$ is proper.  Since $G$ is Keller, it is
a local biholomorphism.  Its image in $U$ is open, and properness makes that
image closed.  The neighborhood $U$ intersects the Zariski-open dense set of
targets having generic fiber cardinality $d$, so the image is nonempty.
Connectedness of $U$ therefore makes the map surjective.  A proper surjective
local biholomorphism is a finite covering map, so its fiber cardinality is
constant on $U$.  Since that cardinality equals $d$ at a generic target, it
equals $d$ over $b$.
\end{proof}

\subsection{The finite-multiple-root component}

Let $D_h$ be the Zariski closure of the targets for which
\eqref{eq:omega} has a finite multiple root.  Since
\[
 \partial_s\Omega_{p,q,r}(0)=2,
\]
such a root is never zero.  Solving $\Omega=\partial_s\Omega=0$ for $q$ and
$r$ gives
\begin{equation}
 \boxed{
 \begin{aligned}
 \Psi_h:\A^1_p\times\mathbb{G}_{m,s}&\longrightarrow\A^3_{p,q,r},\\
 (p,s)&\longmapsto
 \left(
 p,
 p^6h(p)s^4+3ps+\frac1s,
 s-ps^3-\frac23p^6h(p)s^6
 \right).
 \end{aligned}}
 \label{eq:parametrization}
\end{equation}

\begin{lemma}
\label{lem:irreducible}
The variety $D_h$ is an irreducible hypersurface.  It dominates the $p$-line,
is distinct from every vertical hyperplane $V(p-\alpha)$, and contains no
vertical hyperplane.
\end{lemma}

\begin{proof}
The domain $\A^1\times\Gm$ of $\Psi_h$ is irreducible, so the closure of its
image is irreducible.  Its dimension is two: the first target coordinate is
$p$, and
\[
 \frac{\partial q}{\partial s}
 =4p^6h(p)s^3+3p-\frac1{s^2}
\]
is not identically zero.  Hence the differential of $\Psi_h$ has rank two on
a nonempty Zariski-open subset.  By the characteristic-zero dimension theorem
for images, the closure $D_h$ has dimension two.  It is therefore an
irreducible surface in $\A^3$, hence an irreducible hypersurface.

The first coordinate of $\Psi_h(p,s)$ is the freely varying parameter $p$, so
$D_h$ dominates the $p$-line.  A vertical hyperplane does not dominate the
$p$-line, and therefore cannot equal $D_h$.  Finally, if a vertical
hyperplane $H_\alpha=V(p-\alpha)$ were contained in $D_h$, then the two
irreducible closed sets $H_\alpha$ and $D_h$ would have the same dimension
(two).  A proper inclusion of irreducible closed subsets strictly lowers
dimension, so the inclusion would force $H_\alpha=D_h$, contradicting the
preceding paragraph.  Thus no vertical hyperplane is contained in $D_h$.
\end{proof}

\subsection{The reduced nonproperness theorem}

\begin{theorem}
\label{thm:nonproperness}
For every nonzero $h\in\C[T]$,
\begin{equation}
 \boxed{
 S_{F_h}=D_h\cup V(a_h)
 =D_h\cup V\bigl(p^6h(p)\bigr)}
 \label{eq:nonproperness}
\end{equation}
as reduced algebraic sets.
\end{theorem}

\begin{proof}
We prove both inclusions.

\smallskip
\emph{First inclusion: $S_{F_h}\subseteq D_h\cup V(a_h)$.}
Take $b_0=(p_0,q_0,r_0)\notin D_h\cup V(a_h)$.  Then
$a_h(p_0)\ne0$, and $\Omega_{b_0}$ is a degree-six polynomial with six
simple finite roots.  Choose a sufficiently small compact target neighborhood
$K$ of $b_0$ disjoint from $D_h\cup V(a_h)$.  On $K$, the leading coefficient
$a_h(p)$ stays bounded away from zero.  The coefficients of $\Omega$ remain
bounded, so the usual polynomial root bound places every root in one fixed
compact disk.

The incidence set
\[
 \mathcal I_K=
 \{(b,s)\in K\times\C:\Omega_b(s)=0\}
\]
is therefore compact.  No point of $\mathcal I_K$ satisfies
$\partial_s\Omega=0$, because $K\cap D_h=\varnothing$.  Hence, by
\eqref{eq:derivative}, $|D|$ has a positive minimum on $\mathcal I_K$.
Formula \eqref{eq:reconstruction} then shows that all corresponding source
points are uniformly bounded.

There are no additional source points on $A=0$: that equation forces $P=0$,
whereas $p$ stays away from zero on $K$.  Consequently $F_h^{-1}(K)$ is
compact, and $b_0\notin S_{F_h}$.  This proves the first inclusion.

\smallskip
\emph{Second inclusion, finite component: $D_h\subseteq S_{F_h}$.}
Let $(p,q,r)=\Psi_h(p,t)$, where $t\ne0$.  For $\varepsilon\ne0$, set
\[
 y_\varepsilon=\frac{1-\varepsilon}{t},
 \qquad
 x_\varepsilon=\frac{t}{\varepsilon},
 \qquad
 z_\varepsilon
 =p\varepsilon^3-y_\varepsilon^2(4-ty_\varepsilon)\varepsilon.
\]
Then $A=\varepsilon^{-1}$ and $s=t$.  Direct substitution gives
\[
 F_h(x_\varepsilon,y_\varepsilon,z_\varepsilon)
 =\left(p,q-\frac{\varepsilon}{t},r+\varepsilon t\right)
 \longrightarrow(p,q,r).
\]
At the same time, $|x_\varepsilon|=|t/\varepsilon|\to\infty$.  Thus every
point in the image of $\Psi_h$ belongs to $S_{F_h}$.  The image of $\Psi_h$
is constructible and Zariski dense in $D_h$; hence it contains a Zariski-open
dense subset of $D_h$, which is also Euclidean dense.  Lemma
\ref{lem:closed} now gives $D_h\subseteq S_{F_h}$.

\smallskip
\emph{Second inclusion, vertical components: $V(a_h)\subseteq S_{F_h}$.}
The distinct roots of $a_h(p)=\frac13p^6h(p)$ are $0$ together with the
distinct nonzero roots of $h$.  Fix one such root $\alpha$, and consider a
generic target on $H_\alpha=V(p-\alpha)$ outside $D_h$.

If $\alpha\ne0$, then $h(\alpha)=0$, and \eqref{eq:omega} specializes to
\[
 2\alpha s^3-qs^2+2s-r=0.
\]
For a target outside $D_h$, this cubic has three simple roots.
Equation \eqref{eq:reconstruction} gives exactly three source points with
$A\ne0$.  There are no source points with $A=0$, because $A=0$ forces
$P=0\ne\alpha$.  Hence the fiber has three points.

If $\alpha=0$, then, for generic $q\ne0$, equation \eqref{eq:omega} becomes
\[
 -qs^2+2s-r=0.
\]
Outside $D_h$, it has two simple finite roots and hence gives two source
points with $A\ne0$.  There is exactly one further source point on $A=0$.
Indeed, $A=0$ implies $xy=-1$, $B=y^2$, and $P=0$.  Moreover,
\[
 Q_h=-2y,
 \qquad
 R_h=5x-x^3z-\frac23y^6h(0).
\]
Thus, when $q\ne0$,
\[
 y=-\frac q2,
 \qquad
 x=\frac2q,
 \qquad
 z=\frac{5x-\frac23y^6h(0)-r}{x^3},
\]
which is one and only one point.  The generic fiber over $H_0$ therefore also
has three points.

By Lemma \ref{lem:keller-degree}, $F_h$ has generic degree six.  Lemma
\ref{lem:fiber-loss} now implies that every target in these generic open
subsets of the vertical hyperplanes lies in $S_{F_h}$, because its fiber has
only three points.  Each such open subset is Euclidean dense in its
hyperplane: Lemma \ref{lem:irreducible} ensures that $D_h$ contains no
vertical hyperplane.  Since $S_{F_h}$ is closed by Lemma \ref{lem:closed},
every vertical hyperplane $H_\alpha$ is contained in $S_{F_h}$.

This proves $V(a_h)\subseteq S_{F_h}$ and completes
\eqref{eq:nonproperness}.
\end{proof}


\subsection{Earlier pair-specific explicit asymptotic witnesses}
\label{subsec:explicit-witnesses}

An earlier accepted checkpoint established the two concrete vertical loci by
explicit escaping inverse-root branches.  Those witnesses remain useful
independent checks, but the final universal theorem uses the fiber-loss route:
local proper Keller coverings have locally constant finite fibers, dense good
targets have exactly six simple inverse roots, and every leading-locus fiber
has cardinality at most three.  The explicit pair witnesses are recorded here
because they give a second proof for the two original examples.

For the common component $p=0$, let $\varepsilon\to0$, put
\[
 s_\varepsilon=\varepsilon^{-5},
 \qquad
 d(q)=-\frac{q}{10},
 \qquad
 w_\varepsilon=c+d(q)\varepsilon^2,
 \qquad
 p_\varepsilon=\varepsilon^3w_\varepsilon,
\]
and choose
\[
 c^5=-6 \quad\text{for CEX-004},
 \qquad
 c^5=4 \quad\text{for CEX-006}.
\]
For a prescribed target $(0,q,r)$ define
\[
 q_\varepsilon
 =a_h(p_\varepsilon)s_\varepsilon^4
   +2p_\varepsilon s_\varepsilon
   +\frac{2}{s_\varepsilon}
   -\frac{r}{s_\varepsilon^2}.
\]
Then $s_\varepsilon$ is an exact root of
$\Omega_{p_\varepsilon,q_\varepsilon,r}$.  The cancellation conditions on
$c$ give
\[
 p_\varepsilon\longrightarrow0,
 \qquad
 q_\varepsilon\longrightarrow q.
\]
The evaluated derivative is eventually nonzero, so the checked reconstruction
map produces exact source points mapping to
$(p_\varepsilon,q_\varepsilon,r)$.  Moreover the reconstructed second source
coordinate satisfies
\[
 y_\varepsilon\varepsilon^2\longrightarrow3c\ne0,
\]
and hence the sources escape.  This proves the entire hyperplane $p=0$ is
contained in each pair-specific nonproperness set.  It corrects the tempting
but insufficient direct $A=0$ family, which approaches only the line
$(p,q)=(0,0)$.

For the additional CEX-004 hyperplane, set $\alpha=-\frac14$ and let
$s\to\infty$.  For a prescribed target $(\alpha,q,r)$ define
\[
 p_s=\alpha+\frac{1536}{s^3}+\frac{3072q}{s^4},
 \qquad
 q_s=a_{004}(p_s)s^4+2p_ss+\frac2s-\frac{r}{s^2}.
\]
Again $s$ is an exact root of $\Omega_{p_s,q_s,r}$, and the checked limits are
\[
 p_s\longrightarrow-\frac14,
 \qquad
 q_s\longrightarrow q,
 \qquad
 s^{-2}\Omega'_{p_s,q_s,r}(s)\longrightarrow\frac32,
 \qquad
 s^{-1}y_s\longrightarrow-\frac34.
\]
Thus the root is eventually simple, reconstruction is exact, and the source
sequence escapes.  Consequently the whole hyperplane $p=-\frac14$ lies in
$S_{004}$.

The formal forward inclusion uses the inverse chart rather than a compact
neighborhood argument.  If $F_h(u_n)\to b$ and $a_h(b_0)\ne0$, then eventually
$A(u_n)\ne0$ and the chart roots
$s_n=x_n/A(u_n)$ satisfy the inverse equation.  A direct Cauchy-type estimate
bounds $s_n$; after passage to a convergent subsequence, a nonzero limiting
derivative would make the reconstructed sources converge, contradicting
escape.  The limiting root is therefore multiple and nonzero, so $b$ lies in
the finite critical image, hence in $D_h$.  Lean proves the universal inclusion
\[
 \operatorname{NonpropernessSet}(F_h)
 \subseteq D_h\cup V(a_h)
\]
for every deformation polynomial $h$.

To pass from the directly parametrized critical image to its Zariski closure
for the two examples, Lean does not infer Euclidean closure from Zariski
closure.  Instead it verifies exact elimination certificates.  Two Euclidean
remainder identities reduce $\Omega$ and $\Omega'$ to a cubic $H_3$ and a
quadratic $H_2$, and the fixed $5\times5$ Sylvester determinant gives
\[
 \operatorname{Res}(H_3,H_2)
 =27p^{10}\delta_{004}(p,q,r)
\]
for CEX-004 and
\[
 \operatorname{Res}(H_3,H_2)
 =-\frac{243}{8}p^{10}\delta_{006}(p,q,r)
\]
for CEX-006.  Off the leading locus, vanishing of the pair-specific factor
produces a common finite root of $\Omega$ and $\Omega'$, hence a point of the
critical image.  On the leading locus, the explicit vertical branches above
already give nonproperness.  This closes the full reverse inclusions
$D_{004}\subseteq S_{004}$ and $D_{006}\subseteq S_{006}$ inside Lean.

\subsection{Component count}

\begin{corollary}[General component formula]
\label{cor:count}
For every nonzero $h\in\C[T]$,
\begin{equation}
 \boxed{
 \#\Irr(S_{F_h})
 =2+\#\{\alpha\in\C^\times:h(\alpha)=0\}.}
 \label{eq:count}
\end{equation}
The roots are counted without multiplicity.
\end{corollary}

\begin{proof}
As a reduced set,
\[
 V(p^6h(p))
 =V(p)\cup
 \bigcup_{\substack{\alpha\in\C^\times\\h(\alpha)=0}}V(p-\alpha).
\]
These vertical hyperplanes are irreducible and pairwise distinct.  Lemma
\ref{lem:irreducible} says that $D_h$ is irreducible, is not contained in any
vertical hyperplane, and contains no vertical hyperplane.  Hence the displayed
finite union is irredundant.  Theorem \ref{thm:nonproperness} therefore
decomposes $S_{F_h}$ into exactly one finite-multiple-root component, one
component $V(p)$, and one component for every distinct nonzero root of $h$.
\end{proof}

\subsection{CEX-004 versus CEX-006}

For CEX-004,
\[
 h_{004}(T)=1+4T.
\]
It has the single nonzero root $-\frac14$.  Hence
\[
 S_{004}=D_{004}\cup V(p)\cup V(1+4p),
 \qquad
 \#\Irr(S_{004})=3.
\]
For CEX-006,
\[
 h_{006}(T)=-\frac32.
\]
It has no roots.  Hence
\[
 S_{006}=D_{006}\cup V(p),
 \qquad
 \#\Irr(S_{006})=2.
\]

\begin{lemma}[Left-right invariance]
\label{lem:invariance}
If $G=L\circ F\circ R$, where $L$ and $R$ are polynomial automorphisms of
$\C^n$, then
\[
 S_G=L(S_F).
\]
\end{lemma}

\begin{proof}
A polynomial automorphism is proper because its inverse is continuous and
maps compact sets to compact sets.  Consequently a sequence $u_k$ escapes to
infinity if and only if $R(u_k)$ escapes to infinity.  Also,
\[
 G(u_k)\to b
 \quad\Longleftrightarrow\quad
 F(R(u_k))\to L^{-1}(b).
\]
The defining sequence condition for nonproperness therefore gives
$b\in S_G$ if and only if $L^{-1}(b)\in S_F$.
\end{proof}

An algebraic automorphism maps irreducible components bijectively.  Thus
left-right equivalent maps have nonproperness sets with the same number of
irreducible components.  The counts $3$ and $2$ are different, so CEX-004
and CEX-006 are not left-right polynomially equivalent.  This proves
Theorem A.

\subsection{The infinite family}

For $m\ge0$, define
\[
 h_m(T)=\prod_{i=1}^{m}(T-i),
\]
where the empty product $h_0$ is $1$.  The polynomial $h_m$ has exactly $m$
distinct nonzero roots.  Lemma \ref{lem:keller-degree} gives
\[
 \det JF_{h_m}=-2,
 \qquad
 \mu(F_{h_m})=6.
\]
In particular, $F_{h_m}$ is a Keller map but not a polynomial automorphism,
because an automorphism has generic degree one.  Thus every $F_{h_m}$ is a
Keller counterexample.  Corollary \ref{cor:count} gives
\[
 \#\Irr(S_{F_{h_m}})=m+2.
\]
If $m\ne n$, these component counts differ, so Lemma \ref{lem:invariance}
rules out a left-right polynomial equivalence between $F_{h_m}$ and
$F_{h_n}$.  The family $\{F_{h_m}\}_{m\ge0}$ is therefore infinite, pairwise
inequivalent, and has constant generic degree six.  This proves Theorem B.

\section{Lean 4 formalization}

\subsection{Pinned environment and reproducible audit}

The project is pinned by
\[
 \texttt{leanprover/lean4:v4.32.2}
\]
and mathlib commit
\[
 \texttt{905b95818eb32af7874a58b427f50c1711a5e96c}.
\]
The PowerShell entry point \texttt{build\_and\_audit.ps1} verifies the exact
toolchain pin, resolves that pinned Lake executable directly, builds the
umbrella and axiom-audit targets, scans the project Lean sources for
\texttt{axiom}, \texttt{sorry}, or \texttt{admit}, parses every
\texttt{\#print axioms} report, and fails on missing, duplicate, unexpected,
or non-whitelisted results.  It performs no dependency update or cache
download.

The immutable CP-07 replay reported
\begin{lstlisting}[language={}]
Build completed successfully (3507 jobs).
AXIOM_WHITELIST_PASSED (92 reports)
BUILD_AND_AUDIT_PASSED
\end{lstlisting}
Every audited endpoint reported only
\[
 \texttt{propext},\qquad
 \texttt{Classical.choice},\qquad
 \texttt{Quot.sound}.
\]
The CP-07 formal-input manifest contains 53 verified entries.  After the final
paper and documentation synchronization, the release manifest is regenerated
and reverified before packaging.

\subsection{Module map}

The root module has 32 direct runtime imports.  There are 48 submodules under
\texttt{DegreeSixKeller/}: 33 non-Spike modules and 15 Spike-named modules.
The separate \texttt{AxiomAudit} imports the umbrella and is compiled as a
second target.  Two historically named support modules,
\texttt{FunctionFieldSpike} and \texttt{GenericDegreeGeneratorSpike}, remain
live dependencies of the generic-degree layer; the other 13 Spike modules are
outside the production closure and retained as proof-development evidence.

\begin{longtable}{>{\raggedright\arraybackslash\ttfamily}p{0.34\textwidth}>{\raggedright\arraybackslash}p{0.60\textwidth}}
\toprule
Module & Kernel-checked content \\
\midrule
\endfirsthead
\toprule
Module & Kernel-checked content \\
\midrule
\endhead
FetaCore, InverseChart & Point maps, inverse polynomial, exact chart,
derivative, reconstruction, and leading-coefficient locus. \\
PolynomialModel, JacobianKeller & Polynomial coordinate triples, exact
evaluation bridge, formal Jacobian, determinant $-2$, and Keller certificate. \\
FunctionFieldSpike, GenericDegreeGenerator\-Spike, GenericDegree & Source and
coordinate function fields, inverse generator, rational-function transport,
and genuine finrank-six theorem. \\
PolynomialAutomorphism, GenericDegreeAutomorphism & Polynomial inverse data,
left-right bridge, automorphism degree one, and nonautomorphism of every
nonzero family member. \\
AffineZariski, ComponentCounting & Affine-Zariski topology and exact component
counting through mathlib's \texttt{irreducibleComponents}. \\
CriticalImageGeometry, GeneralFiniteGeometry, FiniteComponentDimension &
Universal critical image, exact eliminant zero locus, irreducibility,
hypersurface, dominance, nonverticality, and topological Krull dimension two. \\
KellerLocalHomeomorph, FiberLoss & Keller local homeomorphisms, compact local
restrictions, proper coverings, monodromy, and locally constant fiber size. \\
GenericFiberSix, GoodTargetDensity & Exact six-point fibers on a dense set of
simple-root targets. \\
LeadingFiberBounds & Finiteness and \texttt{encard} at most three on the
complete leading locus, including $p=0$. \\
AsymptoticExclusion, GeneralNonproperness & Universal forward inclusion,
fiber-loss reverse inclusion, and exact reduced nonproperness equality. \\
GeneralComponents & Distinct nonzero-root indexing, edge cases, exact
irreducible-component Finsets, and the formula
$2+\#\{\alpha\ne0:h(\alpha)=0\}$. \\
InfiniteFamily & Product parameters, exact roots and counts, Keller/degree-six
certificates, infinite range, and pairwise polynomial inequivalence. \\
Pair-specific legacy modules & Independent explicit escape/resultant proofs
for CEX-004 and CEX-006 and their original component-count contradiction. \\
MainTheorems & Parameter-free Lean declarations \texttt{theoremA} and
\texttt{theoremB}. \\
AxiomAudit & The 92 public \texttt{\#print axioms} probes used by the enforced
trust audit. \\
\bottomrule
\end{longtable}

\subsection{Exact checked endpoint declarations}

Representative completed declarations are:
\begin{lstlisting}[language={}]
theorem Fh_jacobianDet (h : Complex[X]) :
    jacobianDet (FhPolynomial h) = MvPolynomial.C (-2)

theorem Fh_genericDegree_six (h : Complex[X]) (hh : Not (h = 0)) :
    genericDegree (FhPolynomial h) = 6

theorem Fh_reducedNonpropernessEquality
    (h : Complex[X]) (hh : Not (h = 0)) :
    NonpropernessSet (Fh h) = generalReducedCandidate h

theorem Fh_componentCount (h : Complex[X]) (hh : Not (h = 0)) :
    algebraicComponentCount (NonpropernessSet (Fh h)) =
      2 + (nonzeroRoots h).card

theorem hFamily_componentCount (m : Nat) :
    algebraicComponentCount
      (NonpropernessSet (Fh (hFamily m))) = m + 2

#check DegreeSixKeller.theoremA
#check DegreeSixKeller.theoremB
\end{lstlisting}
Here \texttt{generalReducedCandidate h} expands to the finite component union
the zero locus of $p^6h(p)$.  The pair specializations recover
\[
 \operatorname{NonpropernessSet}(F_{004})
 =\operatorname{finiteComponent}(\eta_{004})
   \cup\operatorname{pHyperplane}(0)
   \cup\operatorname{pHyperplane}\!\left(-\frac14\right)
\]
and
\[
 \operatorname{NonpropernessSet}(F_{006})
 =\operatorname{finiteComponent}(\eta_{006})
   \cup\operatorname{pHyperplane}(0).
\]
The declaration \texttt{theoremA} additionally packages the two Keller,
generic-degree-six, and nonautomorphism certificates.  The declaration
\texttt{theoremB} quantifies over every $m$, packages count $m+2$, and proves
both literal infinitude of the range and pairwise polynomial left-right
inequivalence on that range.  Neither declaration has a theorem parameter or
an unproved geometric hypothesis.

\subsection{From pair closure to the universal theorem}

The two irreducibility obligations are closed by treating the critical
parametrization as a regular map on the principal open $D(s)$ in affine
$2$-space.  That domain is a nonempty open subset of irreducible affine space;
its continuous image is irreducible, and the closure remains irreducible.

The two nonverticality obligations are closed by the exact elimination factors
$\delta_{004}$ and $\delta_{006}$.  Each finite component lies in the
corresponding zero locus.  On every slice $p=\alpha$, $r=0$, the resulting
univariate polynomial in $q$ has a fixed nonzero quadratic coefficient:
$-243$ for CEX-004 and $216$ for CEX-006.  Hence every vertical hyperplane
contains a point outside that zero locus and cannot lie inside the finite
component.

The earlier two reduced-set equalities are closed by combining:
\begin{enumerate}
 \item explicit escaping sequences for the critical image and all pair-specific
       vertical components;
 \item the exact reduced-resultant bridge from the finite Zariski closure back
       to a common inverse root off the leading locus; and
 \item the asymptotic-exclusion theorem proving that no other escaping limit is
       possible.
\end{enumerate}
This is a direct formal proof of both inclusions, not a theorem parameter or an
imported geometric certificate.

The final family-wide reverse inclusion uses the accepted fiber-loss argument,
not an unavailable blanket identification of arbitrary fibers with
function-field finrank.  Away from the nonproperness set, local properness and
the Keller local-homeomorphism theorem produce a finite covering over a target
ball.  Monodromy makes its fiber cardinality constant.  Density supplies a
good target in that ball with six simple inverse roots, so the original fiber
also has six points.  On the leading locus, the algebraic analysis proves
finiteness and \texttt{encard} at most three; hence those targets must be
nonproper.  This covers $p=0$, nonzero roots, repeated roots, and constants.

\subsection{Human-to-Lean correspondence}

\begin{center}
\begin{tabularx}{\textwidth}{>{\raggedright\arraybackslash\bfseries}p{0.36\textwidth}X}
\toprule
Human proof step & Formal status \\
\midrule
Coordinate definitions and inverse polynomial & Checked exactly. \\
Polynomial coordinate model and exact evaluation bridge & Checked for every
$h$. \\
Jacobian determinant $-2$ & Checked for every $h$. \\
Generic degree six & Checked for every nonzero $h$ using a genuine
function-field finrank. \\
Polynomial automorphisms have generic degree one & Checked; consequently every
nonzero family member is nonautomorphic. \\
Closedness of $S_G$ & Checked for the sequence-defined
\texttt{NonpropernessSet}. \\
Fiber-loss theorem & Checked through local proper Keller coverings, monodromy,
dense good targets, and exact six-point fibers. \\
Inverse chart, derivative, and reconstruction & Checked exactly, including
exact mapping from every finite simple root and reconstruction of every source
point on $A\ne0$. \\
Critical-image irreducibility and closure irreducibility & Checked for every
deformation polynomial $h$. \\
Critical image contained in nonproperness & Checked by the explicit
$\varepsilon$-escape family. \\
Complete leading locus & Checked by finite fiber bounds with
\texttt{encard} at most three, including the entire $p=0$ hyperplane. \\
$S_{F_h}\subseteq D_h\cup V(a_h)$ & Checked universally for every $h$ by root
bounds, subsequence extraction, and reconstruction contradiction. \\
$D_h\cup V(a_h)\subseteq S_{F_h}$ & Checked for every nonzero $h$ by the
fiber-loss route. \\
Finite-component geometry & Closed, irreducible, exact hypersurface,
topological Krull dimension two, dominant, and contains no vertical
hyperplane. \\
Universal irreducible-component formula & Checked using actual
\texttt{irreducibleComponents}. \\
Left-right transport of nonproperness & Checked for homeomorphisms with
explicit escape preservation. \\
Component-count invariance & Checked for a target map equipped with a
compatible affine-Zariski homeomorphism. \\
Polynomial automorphisms supply those data & Checked through polynomial
forward/inverse coordinate data and the induced Euclidean and affine-Zariski
homeomorphisms. \\
Final $3\ne2$ contradiction in the formal interface & Checked
unconditionally. \\
Explicit product family and count $m+2$ & Checked, including the empty product,
exact distinct roots, infinite range, and pairwise inequivalence. \\
Parameter-free Theorems A and B & Checked unconditionally. \\
\bottomrule
\end{tabularx}
\end{center}

\subsection{Axiom, hole, and integrity audit}

The audited public declarations report only
\[
 \texttt{propext},\qquad
 \texttt{Classical.choice},\qquad
 \texttt{Quot.sound}.
\]
These are ordinary Lean/mathlib foundational dependencies, not
project-specific mathematical assumptions.  The use of
\texttt{Classical.choice} is visible in constructions such as chosen fifth
roots, chosen principal-open pullback data, and noncomputable component
counts.  The source scan found no project-local \texttt{axiom},
\texttt{sorry}, or \texttt{admit}.  No CAS output, floating-point computation,
or numerical resultant is trusted: the polynomial and determinant identities
are elaborated and checked by Lean.

The final \texttt{SOURCE\_MANIFEST.sha256} is generated after the paper and
documentation edits and verified entry by entry before packaging.  This
integrity replay is independent of the logical kernel audit.

\subsection{Exact remaining formal boundary}

The mathematical content of Theorems A and B is closed in Lean.  The formal
equalities concern underlying reduced sets, and the component invariant is the
cardinality of mathlib's actual irreducible-component set.  The only stated
boundary is scheme-theoretic: multiplicities, valuations, normalizations, and
scheme structures on the nonproperness locus are not defined or claimed.

\section{Dual closure audit}

\subsection{Mathematical closure audit: passed at H1}

The human argument closes every stated mathematical obligation:
\begin{itemize}[leftmargin=2.1em,itemsep=0.25em]
 \item The exact reduced-set equality, component formula, present-pair
       polynomial inequivalence, and infinite degree-six theorem are proved.
 \item The hypothesis $h\ne0$ is used exactly where needed to keep the generic
       degree equal to six; distinct roots are counted without multiplicity.
 \item Constant $h$, a root at $0$, repeated roots, the special component
       $p=0$, and the empty product $m=0$ are covered.
 \item Nonproperness is established independently of component counting, so
       the decomposition and the invariant are not used circularly.
 \item Both possible containments between the finite component and a vertical
       hyperplane are excluded.
 \item Equality always means equality of underlying reduced algebraic sets.
       No scheme structure is tacitly imported.
\end{itemize}

The completed Lean artifact independently checks the entire universal spine:
polynomial and point-map agreement, Jacobian and degree, nonautomorphism,
inverse formulas, finite geometry, proper-covering fiber loss, the reduced-set
equality, actual component formula, and both parameter-free main theorems.

\subsection{Formal closure audit: Theorems A and B passed}

The parameter-free \texttt{theoremA} and \texttt{theoremB}, together with all
their public supporting endpoints, are accepted by the pinned kernel.  The
3,507-job replay and 92-report enforced axiom audit supply the immutable
formal-closure evidence.

\section{Remaining frontier}

There is no remaining theorem-level action item for the claims in this paper.
An optional future project could define and study a scheme-theoretic or
valuation-theoretic refinement of the nonproperness locus.  Such a refinement
would be a new theorem with new invariance obligations; factors and
multiplicities in a raw resultant are not promoted to invariants here.

\begin{thebibliography}{1}

\bibitem{source}
\emph{A Counterexample to the Jacobian Conjecture}, working manuscript,
20 July 2026, Theorem 5.2.

\end{thebibliography}

\end{document}
